\documentclass{article}
\usepackage{amsmath,amssymb}
\begin{document}

\section*{Lean-friendly proof request: Higham Ch. 6 Lemma 6.6 real rectangular variant}

I am formalizing a matrix norm result from Higham Chapter 6 in Lean 4 / Mathlib.  Please give a rigorous, Lean-friendly proof plan, with intermediate lemmas and proof details, for bridging an existing complex SVD/operator-norm theorem back to an older real rectangular operator-bound API.

\subsection*{Local real API}

For a real rectangular matrix \(A : \mathrm{Fin}\ m \to \mathrm{Fin}\ n \to \mathbb{R}\), the repository has:
\[
  \operatorname{absMatrixRect}(A)_{ij} := |A_{ij}|.
\]
It has a real Euclidean vector norm \(\operatorname{vecNorm2}\) and rectangular multiplication \(\operatorname{rectMatMulVec}\).  The older real rectangular operator-bound predicate is:
\[
  \operatorname{rectOpNorm2Le}(M,c)
  \quad:\!\!\iff\quad
  \forall x : \mathrm{Fin}\ n \to \mathbb{R},\
  \|\operatorname{rectMatMulVec}(M,x)\|_2 \le c\,\|x\|_2.
\]
Already proved in this real API:
\[
  \operatorname{rectOpNorm2Le}(|A|,\|A\|_F).
\]
This only gives a Frobenius cap.  The desired remaining variant is rank-sensitive.

\subsection*{Local complex API already proved}

For a complex matrix \(B : \mathrm{Fin}\ m \to \mathrm{Fin}\ n \to \mathbb{C}\), the repository has:
\[
  \operatorname{complexMatrixOp2}(B)
\]
as the Euclidean operator norm, \(\operatorname{complexAbsMatrix}(B)_{ij} = |B_{ij}|\) embedded in \(\mathbb{C}\), and
\[
  \operatorname{complexMatrixRank}(B) := \operatorname{Matrix.rank}(B).
\]
The following theorem is already proved for \(0<n\):
\[
  \operatorname{complexMatrixOp2}(|B|)
  \le
  \sqrt{\operatorname{rank}(B)}\,\operatorname{complexMatrixOp2}(B).
\]
There are also concrete \(p=2\) wrappers, but the clean route for this request is probably through \(\operatorname{complexMatrixOp2}\).

\subsection*{Likely target theorem}

Define the complexification of a real rectangular matrix by
\[
  A_\mathbb{C}(i,j) := (A_{ij} : \mathbb{C}),
\]
and define
\[
  \operatorname{realRectMatrixRank}(A)
  := \operatorname{Matrix.rank}(A_\mathbb{C}).
\]
I want a theorem of the following form:
\[
\begin{aligned}
&0<n,\quad 0\le c,\quad \operatorname{rectOpNorm2Le}(A,c)\\
&\qquad\Longrightarrow
\operatorname{rectOpNorm2Le}\bigl(|A|,\sqrt{\operatorname{realRectMatrixRank}(A)}\,c\bigr).
\end{aligned}
\]

\subsection*{What I need from you}

Please give a detailed proof plan suitable for Lean formalization.  In particular:

\begin{enumerate}
\item Prove, mathematically and with Lean-style helper statements, that a real rectangular bound on \(A\) implies the complex Euclidean operator norm bound
\[
  \operatorname{complexMatrixOp2}(A_\mathbb{C}) \le c.
\]
The key issue is that the real predicate only quantifies over real vectors, whereas the complex operator norm quantifies over complex vectors.  I expect the proof to split \(z=x+iy\) and use
\[
  \|A_\mathbb{C}z\|_2^2 = \|Ax\|_2^2+\|Ay\|_2^2
  \le c^2(\|x\|_2^2+\|y\|_2^2).
\]

\item Prove the reverse bridge needed at the end: if
\[
  \operatorname{complexMatrixOp2}((|A|)_\mathbb{C}) \le C
\]
and \(0\le C\), then
\[
  \operatorname{rectOpNorm2Le}(|A|,C).
\]
This should follow by embedding a real vector \(x\) as a complex vector and comparing real and complex Euclidean norms.

\item Prove the compatibility facts:
\[
  (|A|)_\mathbb{C} = |A_\mathbb{C}|,
  \qquad
  \operatorname{complexMatrixRank}(A_\mathbb{C})
  = \operatorname{realRectMatrixRank}(A)
\]
where the second is probably by definition if we define real rank this way.

\item Combine these with the existing complex theorem to get the target bound.  Please spell out all scalar monotonicity steps, especially why
\[
  \sqrt{r}\operatorname{complexMatrixOp2}(A_\mathbb{C})
  \le
  \sqrt{r}\,c
\]
is valid and why the final radius is nonnegative.

\item Warn about any hidden assumptions.  For example: do we need \(0<n\), \(0<m\), \(0\le c\), or can some be derived?  Is defining real rank as complex rank mathematically acceptable for this bridge?
\end{enumerate}

Please do not rely on project files or book text.  I need a rigorous mathematical proof route and Lean-friendly lemma decomposition, not a polished essay.

\end{document}
